% =====================================================================
%  setup/preamble.tex -- packages and layout for the SEF639 Final Report
%
%  You should not need to edit this file. Your own details belong in
%  main.tex. If you add packages, add them at the end of main.tex's
%  preamble (before \begin{document}) so that updates to this template
%  can be copied in without losing your changes.
% =====================================================================
\makeatletter

% ---------------------------------------------------------------------
% Draft or final mode (set \reportmode in main.tex)
% ---------------------------------------------------------------------
\newif\iffinalmode
\def\fyp@finalstr{final}
\ifx\reportmode\fyp@finalstr \finalmodetrue \fi

% ---------------------------------------------------------------------
% Language and fonts (LuaLaTeX)
% ---------------------------------------------------------------------
\usepackage{fontspec}            % Latin Modern by default; see README
\usepackage{polyglossia}
\setdefaultlanguage[variant=british]{english}
\usepackage[autostyle]{csquotes}
\usepackage{microtype}

% ---------------------------------------------------------------------
% Page layout
% ---------------------------------------------------------------------
\usepackage[a4paper, top=25mm, bottom=25mm, left=25mm, right=25mm,
  includehead, includefoot, marginparwidth=20mm, marginparsep=2mm]{geometry}
\usepackage{setspace}
\setstretch{1.15}                % slightly more open than single spacing

% ---------------------------------------------------------------------
% Colours (Solent University brand guidelines)
% ---------------------------------------------------------------------
\usepackage[table]{xcolor}
\definecolor{SolentRed}{RGB}{205,20,35}
\definecolor{SkyBlue}{RGB}{1,195,222}
\definecolor{SeaBlue}{RGB}{15,67,95}
% risk ratings (risk register)
\definecolor{RiskLow}{RGB}{198,239,206}
\definecolor{RiskMedium}{RGB}{255,235,156}
\definecolor{RiskHigh}{RGB}{252,192,140}
\definecolor{RiskVeryHigh}{RGB}{255,150,150}

% ---------------------------------------------------------------------
% Figures, tables, lists, units, maths, code
% ---------------------------------------------------------------------
\usepackage{graphicx}
\graphicspath{{figures/}}
\usepackage{booktabs, tabularx, longtable, multirow, array}
\usepackage{enumitem}
\usepackage[font=small, labelfont=bf]{caption}
\usepackage{subcaption}
\usepackage{amsmath, amssymb}
\usepackage{siunitx}
\sisetup{group-separator={,}, group-minimum-digits=4, per-mode=symbol}
\usepackage{listings}
\usepackage{pdfpages}
\usepackage[en-GB]{datetime2}

\usepackage{tikz}
\usetikzlibrary{positioning, arrows.meta, calc, fit, backgrounds,
  shapes.geometric, shapes.misc}
\usepackage{pgfplots}
\pgfplotsset{compat=1.18}
\usepackage{pgfplotstable}
\usepackage{pgfgantt}
\usepackage{tcolorbox}
\tcbuselibrary{skins, breakable}

% ---------------------------------------------------------------------
% Headers and footers: page number bottom centre (Handbook, App. A)
% ---------------------------------------------------------------------
\usepackage[headsepline=false]{scrlayer-scrpage}
\clearpairofpagestyles
\cfoot*{\pagemark}
\pagestyle{scrheadings}

% ---------------------------------------------------------------------
% Draft-mode helpers: todo notes, watermark, header stamp
% ---------------------------------------------------------------------
\iffinalmode
  \usepackage[disable]{todonotes}
\else
  \usepackage[textsize=scriptsize, color=SkyBlue!30, bordercolor=SkyBlue]{todonotes}
  % todo notes are not counted as words
  \NewCommandCopy\fyp@todo\todo
  \renewcommand\todo[2][]{\directlua{fyp_wc.paused = fyp_wc.paused + 1}%
    \fyp@todo[#1]{#2}\directlua{fyp_wc.paused = fyp_wc.paused - 1}}
  \usepackage{draftwatermark}
  \SetWatermarkText{DRAFT}
  \SetWatermarkColor[gray]{0.94}
  \SetWatermarkScale{1.2}
  \ohead*{\scriptsize\textcolor{SolentRed}{\textbf{DRAFT}}\ compiled
    \DTMnow\ \textbar\ main body: \fyp@lastwords\ words}
\fi

% ---------------------------------------------------------------------
% References: Harvard Solent (Solent University's interpretation of
% BS ISO 690), see https://libguides.solent.ac.uk/HarvardSolent
% Same configuration as the SEF639 Project Handbook.
% ---------------------------------------------------------------------
\usepackage[
  backend=biber,
  style=ext-authoryear,
  sorting=nyt,
  maxbibnames=3, minbibnames=1,   % four or more authors -> et al.
  maxcitenames=3, mincitenames=1,
  giveninits=true,
  uniquename=false, uniquelist=false,
  dashed=false,
  date=year,
  urldate=long,
  isbn=false, doi=false, eprint=false
]{biblatex}

% In-text: (Johns 2014) / Johns (2014); several sources separated by ';'
\DeclareDelimFormat{nameyeardelim}{\addspace}
\DeclareDelimFormat{finalnamedelim}{\addspace\bibstring{and}\space}
\DeclareDelimFormat[bib]{finalnamedelim}{\addspace\bibstring{and}\space}
\DeclareDelimFormat{multicitedelim}{\addsemicolon\space}
% Reference list: JOHNS, R., 2014. \emph{Title}. 6th ed. London: Publisher
\DeclareDelimFormat[bib]{nameyeardelim}{\addcomma\space}
\DeclareFieldFormat{biblabeldate}{#1}
\DeclareNameAlias{sortname}{family-given/given-family}
\renewrobustcmd*{\bibinitdelim}{}
\AtBeginBibliography{\renewcommand*{\mkbibnamefamily}[1]{\MakeUppercase{#1}}}
\renewcommand*{\finentrypunct}{}
\DeclareFieldFormat{editortype}{#1}
% ... [viewed 8 October 2026]. Available from: URL
\newcommand*{\solentmonth}[1]{\ifcase#1\or January\or February\or March\or
  April\or May\or June\or July\or August\or September\or October\or
  November\or December\fi}
\DeclareFieldFormat{url}{Available from: \url{#1}}
\renewbibmacro*{url+urldate}{%
  \iffieldundef{url}{}{%
    \iffieldundef{urlyear}{}{%
      \setunit{\addspace}%
      \printtext[brackets]{viewed\space\thefield{urlday}\space
        \solentmonth{\thefield{urlmonth}}\space\thefield{urlyear}}}%
    \setunit{\adddot\space}\printfield{url}}}
% Split the list into 'References' (cited) and 'Bibliography' (consulted
% but not cited -- add those with \nocite{key} in main.tex)
\DeclareBibliographyCategory{cited}
\AtEveryCitekey{\addtocategory{cited}{\thefield{entrykey}}}

% ---------------------------------------------------------------------
% Cross-references and links
% ---------------------------------------------------------------------
\usepackage{xurl}
\urlstyle{same}
% Work-around for a varioref bug (v1.6i, 2026): define the language hook
% before varioref loads. See https://github.com/latex3/latex2e/issues/2112
\providecommand\extrasenglish{}
\usepackage{varioref}
\usepackage{hyperref}
\hypersetup{
  colorlinks=true,
  linkcolor=\iffinalmode SeaBlue\else SolentRed\fi,
  citecolor=SeaBlue,
  urlcolor=SeaBlue,
  breaklinks=true,
  pdfdisplaydoctitle=true,
  bookmarksnumbered=true
}
\usepackage[nameinlink]{cleveref}
\crefname{chapter}{Chapter}{Chapters}
\crefname{section}{Section}{Sections}
\crefname{subsection}{Section}{Sections}
\crefname{appendix}{Appendix}{Appendices}

% ---------------------------------------------------------------------
% Notation, abbreviations and glossary (entries in glossary-entries.tex)
% ---------------------------------------------------------------------
\usepackage[acronym, symbols, nonumberlist, nopostdot, nogroupskip]{glossaries}
\makeglossaries

% ---------------------------------------------------------------------
% Headings required by the Handbook (Appendix A)
% ---------------------------------------------------------------------
\gappto\captionsenglish{%
  \renewcommand\contentsname{CONTENTS}%
  \renewcommand\listfigurename{FIGURES}%
  \renewcommand\listtablename{TABLES}%
  \renewcommand\bibname{References}}
\setkomafont{disposition}{\normalcolor\sffamily\bfseries}
\addtokomafont{chapter}{\color{SeaBlue}}

% =====================================================================
% Your details (set in main.tex)
% =====================================================================
\newcommand\studentname[1]{\gdef\fyp@student{#1}}
\newcommand\studentid[1]{\gdef\fyp@studentid{#1}}
\newcommand\award[1]{\gdef\fyp@award{#1}}
\newcommand\course[1]{\gdef\fyp@course{#1}}
\newcommand\projecttitle[1]{\gdef\fyp@title{#1}}
\newcommand\supervisor[1]{\gdef\fyp@supervisor{#1}}
\newcommand\academicyear[1]{\gdef\fyp@year{#1}}
\newcommand\reportdate[1]{\gdef\fyp@date{#1}}
\newcommand\wordtarget[1]{\gdef\fyp@target{#1}}
\newif\ifwordcountontitle
\studentname{Student Name}\studentid{00000000}\award{BEng (Hons)}
\course{Course Name}\projecttitle{Project Title}\supervisor{Supervisor}
\academicyear{2026--27}\reportdate{\DTMmonthname{\month}\ \the\year}
\wordtarget{10000}

\AtBeginDocument{\hypersetup{pdftitle={\fyp@title},
  pdfauthor={\fyp@student}, pdfsubject={SEF639 Final Report}}}

% =====================================================================
% Main body and word count
% =====================================================================
\directlua{dofile(kpse.find_file("setup/wordcount.lua") or "setup/wordcount.lua")}
\providecommand\fyp@lastwords{--}   % value from the previous run (.aux)

% Headers and footers are built in the output routine: do not count them
\newcommand\fyp@pause{\directlua{fyp_wc.paused = fyp_wc.paused + 1}}
\newcommand\fyp@resume{\directlua{fyp_wc.paused = fyp_wc.paused - 1}}
\AtBeginDocument{\output=\expandafter{\expandafter\fyp@pause\the\output\fyp@resume}}

% Remember chapter titles for the word-count report
\pretocmd{\addchaptertocentry}%
  {\protected@edef\fyp@tmp{#2}\expandafter\global\expandafter\let\csname fyp@chaptitle@#1\endcsname\fyp@tmp}{}{}

\newcommand\startmainmatter{%
  \clearpage\pagenumbering{arabic}% Introduction starts on page 1
  \par\directlua{fyp_wc.active = true}}

\newcommand\stopmainmatter{%
  \par\directlua{fyp_wc.active = false}%
  \edef\fyp@words{\directlua{tex.sprint(fyp_wc.total)}}%
  \immediate\write\@mainaux{\gdef\string\fyp@lastwords{\fyp@words}}%
  \typeout{^^J*** Main body word count: \fyp@words\space
    (target \fyp@target) ***^^J}%
  \ifnum\numexpr10*\fyp@words<\numexpr9*\fyp@target\relax
    \PackageWarningNoLine{SEF639}{Main body has \fyp@words\space words,
      more than 10\% below the target of \fyp@target}\fi
  \ifnum\numexpr10*\fyp@words>\numexpr11*\fyp@target\relax
    \PackageWarningNoLine{SEF639}{Main body has \fyp@words\space words,
      more than 10\% above the target of \fyp@target}\fi}

% Word-count report (draft mode only, printed after the appendices)
\newcommand\FYPwcrow[2]{%
  #1 & \ifcsname fyp@chaptitle@#1\endcsname
         \csname fyp@chaptitle@#1\endcsname\fi & \num{#2}\\}
\newcommand\wordcountreport{%
  \iffinalmode\else
  \clearpage
  \addchap{Word Count Report (draft only)}
  This page is only printed in draft mode. It counts the words of the main
  body as typeset: body text, headings and list items. Captions, tables,
  footnotes, equations, guidance boxes, front matter, references and
  appendices are not counted. The Handbook expects the main body to be
  within $\pm$10\% of the target.
  \begin{center}
  \begin{tabular}{rlr}
    \toprule
    & \textbf{Chapter} & \textbf{Words}\\ \midrule
    \directlua{fyp_wc.report_rows()}
    \midrule
    & \textbf{Main body total} & \textbf{\num{\fyp@words}}\\
    & Target ($\pm$10\%) & \num{\fyp@target}
      (\num{\the\numexpr9*\fyp@target/10\relax}--\num{\the\numexpr11*\fyp@target/10\relax})\\
    \bottomrule
  \end{tabular}
  \end{center}
  \fi}

% =====================================================================
% Front and back matter helpers
% =====================================================================
\newcommand\printnotation{%
  \printglossary[type=symbols, title=NOTATION, toctitle=Notation]%
  \printglossary[type=\acronymtype, title=ABBREVIATIONS, toctitle=Abbreviations]%
  \printglossary[title=GLOSSARY, toctitle=Glossary]}

\newcommand\printreferences{%
  \printbibliography[heading=bibintoc, title=References, category=cited]%
  \printbibliography[heading=bibintoc, title=Bibliography, notcategory=cited]}

% =====================================================================
% Building blocks (see the example chapters for how to use them)
% =====================================================================

% Guidance boxes: shown in draft mode, removed completely in final mode.
% Delete them as you write your own text.
\iffinalmode
  \NewDocumentEnvironment{guidance}{O{Guidance} +b}{}{}
\else
  \NewDocumentEnvironment{guidance}{O{Guidance}}
    {\par\directlua{fyp_wc.paused = fyp_wc.paused + 1}%
     \begin{tcolorbox}[enhanced, breakable, colback=SkyBlue!8,
      colframe=SkyBlue, coltitle=SeaBlue, fonttitle=\sffamily\bfseries\small,
      fontupper=\small, title={#1}, boxrule=0.6pt, arc=1mm,
      before upper={\setlength{\parskip}{4pt}}]}
    {\end{tcolorbox}\par\directlua{fyp_wc.paused = fyp_wc.paused - 1}}
\fi

% Grey placeholder box, e.g. for a missing figure or logo
\newcommand\placeholder[2][\linewidth]{%
  \begin{tcolorbox}[width=#1, colback=black!4, colframe=black!35,
    boxrule=0.5pt, arc=0pt, halign=center, valign=center,
    fontupper=\small\sffamily\color{black!60}, before skip=4pt, after skip=4pt]
  #2\end{tcolorbox}}

% Objectives O1, O2 ... that can be referenced with \cref
\newlist{objectives}{enumerate}{1}
\setlist[objectives]{label=\textbf{O\arabic*}, ref=O\arabic*, leftmargin=*}
\crefname{objectivesi}{objective}{objectives}
\Crefname{objectivesi}{Objective}{Objectives}

% Requirements R1, R2 ... : use \req{req:label} in the first column of a
% (non-tabularx) table; refer to them later with \cref{req:label}
\newcounter{requirement}
\renewcommand\therequirement{R\arabic{requirement}}
\newcommand\req[1]{\refstepcounter{requirement}\label{#1}\textbf{\therequirement}}
\crefname{requirement}{requirement}{requirements}
\Crefname{requirement}{Requirement}{Requirements}

% Test outcomes
\newcommand\pass{\textcolor{green!45!black}{\textbf{Pass}}}
\newcommand\fail{\textcolor{SolentRed}{\textbf{Fail}}}
\newcommand\partpass{\textcolor{orange!80!black}{\textbf{Partial}}}

% Risk register: \riskrating{likelihood 1-5}{severity 1-5}
% prints the product and colours the table cell
\newcommand\riskrating[2]{%
  \edef\fyp@risk{\the\numexpr#1*#2\relax}%
  \ifnum\fyp@risk<5 \cellcolor{RiskLow}\fyp@risk\ (Low)%
  \else\ifnum\fyp@risk<10 \cellcolor{RiskMedium}\fyp@risk\ (Medium)%
  \else\ifnum\fyp@risk<16 \cellcolor{RiskHigh}\fyp@risk\ (High)%
  \else\cellcolor{RiskVeryHigh}\fyp@risk\ (Very high)\fi\fi\fi}

% Weighted decision matrix (use a plain tabular, not tabularx):
%   \resetscores{A,B,C}       at the start of the table
%   \weight{0.3}              prints the weight of the current criterion
%   \score{A}{4}              prints the score and adds weight x score to A
%   \weightsum, \scoretotal{A} in the last row
\newcommand\resetscores[1]{%
  \gdef\fyp@wsum{0}\gdef\fyp@w{0}%
  \@for\fyp@opt:=#1\do{\expandafter\gdef\csname fyp@tot@\fyp@opt\endcsname{0}}}
\newcommand\weight[1]{\gdef\fyp@w{#1}\xdef\fyp@wsum{\fpeval{\fyp@wsum+#1}}#1}
\newcommand\score[2]{#2\expandafter\xdef\csname fyp@tot@#1\endcsname
  {\fpeval{\csname fyp@tot@#1\endcsname+\fyp@w*(#2)}}}
\newcommand\scoretotal[1]{\fpeval{round(\csname fyp@tot@#1\endcsname,2)}}
\newcommand\weightsum{\fpeval{round(\fyp@wsum,2)}%
  \ifdim\fpeval{abs(\fyp@wsum-1)}pt>0.001pt
    \ \textcolor{SolentRed}{(should be 1)}\fi}

% TikZ styles for flowcharts (methodology, system diagrams)
\tikzset{
  fyp node/.style={draw=SeaBlue, thick, align=center, font=\small,
    execute at begin node={\hyphenpenalty=10000},
    minimum height=2.4em, inner sep=4pt},
  fyp process/.style={fyp node, rectangle, fill=SkyBlue!12,
    text width=9em},
  fyp terminal/.style={fyp node, rounded rectangle, fill=SeaBlue!15,
    text width=7em},
  fyp decision/.style={fyp node, diamond, aspect=2.2, fill=SolentRed!8,
    text width=6.5em, inner sep=1pt},
  fyp io/.style={fyp node, trapezium, trapezium left angle=72,
    trapezium right angle=108, trapezium stretches body,
    fill=black!5, text width=8em},
  fyp data/.style={fyp node, cylinder, shape border rotate=90,
    aspect=0.15, fill=black!5, text width=6em},
  fyp deliverable/.style={fyp node, rectangle, double, fill=SkyBlue!25,
    text width=9em},
  fyp arrow/.style={-{Stealth[length=2.5mm]}, thick, SeaBlue},
  fyp group/.style={draw=black!35, dashed, rounded corners, inner sep=8pt,
    fill=black!2}
}

% Code listings
\lstdefinestyle{fyp}{
  basicstyle=\ttfamily\footnotesize, numbers=left,
  numberstyle=\tiny\color{black!50}, frame=single, rulecolor=\color{black!25},
  keywordstyle=\color{SeaBlue}\bfseries, commentstyle=\color{green!40!black},
  stringstyle=\color{SolentRed}, breaklines=true, showstringspaces=false,
  captionpos=b, tabsize=2}
\lstset{style=fyp}

\makeatother
